% ======================================================================
% Reporte MAP: Serie de Fourier Truncada y FFT (Matematica 6)
% Cada integrante escribe SOLO su seccion para evitar conflictos en git.
% Las figuras se generan desde los notebooks en reporte/figuras/
% ======================================================================
\documentclass[11pt,letterpaper]{article}
\usepackage[utf8]{inputenc}
\usepackage[T1]{fontenc}
\usepackage[spanish]{babel}
\usepackage{amsmath,amssymb}
\usepackage{graphicx}
\usepackage{booktabs}
\usepackage[margin=2.5cm]{geometry}
\usepackage{hyperref}
\graphicspath{{figuras/}}

\title{MAP: Serie de Fourier Truncada y FFT}
\author{Christian \and Tony \and Claudia}
\date{Noviembre de 2026}

\begin{document}
\maketitle

\section{Introduccion}
% TODO (Claudia): objetivo del proyecto y estructura del reporte.

\section{Parte 1: Series de Fourier} % Christian
\subsection{Interfaz, coeficientes y serie truncada (inciso 1)}
% TODO
\subsection{Documentacion para $N=3,5,10$ (inciso 2)}
% TODO
\subsection{Respuestas justificadas (inciso 3)}
% TODO: a) convergencia, b) utilidad del ICE, c) iteraciones
\subsection{Funcion generada por IA (incisos 4 a 6)}
% TODO: prompt y respuesta (captura), N=15, grafica

\section{Parte 2: FFT y filtros de imagen} % Tony
\subsection{Implementacion de la FFT de Cooley-Tukey}
% TODO: explicar fft_recursiva, FFT 2D por filas y columnas, pruebas contra NumPy
\subsection{Convencion de frecuencia de corte}
% TODO: radio normalizado r en [0,1] desde el centro; esquina = 1
\subsection{Filtro pasa bajas (inciso 1)}
% TODO: figuras pasa bajas fc = 0.22, 0.15, 0.96 y justificacion
\subsection{Filtro pasa altas (inciso 2)}
% TODO: figuras pasa altas fc = 0.23, 0.11, 0.91 y justificacion
\subsection{Blur Detection (inciso 3)}
% TODO: metodo, umbral, resultados

\section{Parte 2: CNN y comparacion} % Claudia
\subsection{Modelo preentrenado (inciso 4a)}
% TODO
\subsection{Ruido, limpieza con FFT y con CNN (incisos 4b a 4d)}
% TODO
\subsection{Comparacion (inciso 5)}
% TODO

\section{Uso de herramientas de IA}
% TODO: cada consulta con captura del prompt y la respuesta, y como se uso

\section{Conclusiones}
% TODO

\end{document}
